% GENERATED FILE. Edit canonical lessons, then run npm run generate:books.
% canonical-source-hash: fnv1a64:7621daa69cbb19cd
% canonical-lessons: TE-C511-telivitetalu, TE-C511-tikamakapadu, TE-C511-talacuko, TE-C511-nemaruveyu, TE-C511-gurtuceyu

\chapter{Memory and Muddle}
\label{ch:te-memory-and-muddle}
\hlchaptermodality{511}

\begin{chapteropening}
\textbf{By the end of this chapter:} \emph{I can say that I was confused, or that I recalled or reminded someone of something.}

Complete the last lesson of chapter 511: I can say that I was confused, or that I recalled or reminded someone of something.
\end{chapteropening}

\section[\texorpdfstring{\te{తెలివితేటలు} (telivitēṭalu)}{తెలివితేటలు (telivitēṭalu)}]{\te{తెలివితేటలు} (telivitēṭalu) — wits; intelligence}

\label{lesson:TE-C511-telivitetalu}



\begin{quote}
Before the new one: say the Telugu for self-confidence, then the Telugu for curiosity.
\end{quote}

\subsection*{You'll want to know: \te{తెలివితేటలు}}
\textbf{\te{తెలివితేటలు}} — \emph{telivitēṭalu} — \textquotedblleft{}wits; intelligence\textquotedblright{}.

\begin{quote}
\te{ఆమె} \te{తెలివితేటలతో} \te{సమస్యను} \te{పరిష్కరించింది}.
\end{quote}

\begin{quote}
\emph{āme telivitēṭalatō samasyanu pariṣkariñcindi.}
\end{quote}

\begin{quote}
She solved the problem with her wits.
\end{quote}

Plural in form. \te{తెలివి} is understanding.

\begin{cousinweb}
\textbf{\te{తెలివితేటలు}} (\emph{telivitēṭalu}) is native Dravidian: DEDR 3433 lists Telugu \emph{telivi}, \textquotedblleft{}understanding\textquotedblright{}, from \emph{teliyu}, \textquotedblleft{}to be known, to be clear\textquotedblright{}; \emph{tēṭalu} goes with it as a pair word.
\end{cousinweb}

\begin{grammarlens}[title={What you've built}]
One more word for describing what you lived and felt.
\end{grammarlens}

\subsection*{Your turn}
\begin{itemize}
\raggedright
  \item \emph{Say it:} \emph{telivitēṭalu}
  \item \emph{Say it:} \emph{telivitēṭalu}, once more
  \item \emph{Say it:} \emph{āme telivitēṭalatō samasyanu pariṣkariñcindi.} — She solved the problem with her wits.
  \item \emph{Recall:} say the Telugu for self-confidence, then the Telugu for curiosity, then say \emph{telivitēṭalu} again
\end{itemize}

\begin{tcolorbox}[breakable,colback=teal!4,colframe=teal!35!black,title={Before you move on}]
What does \te{తెలివితేటలు} mean? (\textquotedblleft{}Wits; intelligence\textquotedblright{}.) Say it once more.
\end{tcolorbox}
\section[\texorpdfstring{\te{తికమకపడు} (tikamakapaḍu)}{తికమకపడు (tikamakapaḍu)}]{\te{తికమకపడు} (tikamakapaḍu) — to be confused; to be muddled}

\label{lesson:TE-C511-tikamakapadu}



\begin{quote}
Before the new one: say the Telugu for curiosity, then the Telugu for wits.
\end{quote}

\subsection*{You'll want to know: \te{తికమకపడు}}
\textbf{\te{తికమకపడు}} — \emph{tikamakapaḍu} — \textquotedblleft{}to be confused; to be muddled\textquotedblright{}.

\begin{quote}
\te{కొత్త} \te{నగరంలో} \te{దారి} \te{తెలియక} \te{నేను} \te{తికమకపడ్డాను}.
\end{quote}

\begin{quote}
\emph{kotta nagaraṁlō dāri teliyaka nēnu tikamakapaḍḍānu.}
\end{quote}

\begin{quote}
Not knowing the way in the new city, I got muddled.
\end{quote}

Past (I …): \te{తికమకపడ్డాను}, \emph{tikamakapaḍḍānu}.

\te{తికమక}, a muddle: \te{తికమకగా} \te{ఉంది}, it is confusing.

\begin{cousinweb}
\textbf{\te{తికమకపడు}} (\emph{tikamakapaḍu}) is native Dravidian: DEDR 3207 lists Telugu \emph{tikamaka}, \textquotedblleft{}intricacy, confusion\textquotedblright{}, and \emph{tikamakalāḍu}, \textquotedblleft{}to be puzzled\textquotedblright{}, with Tamil \emph{tikai}, \textquotedblleft{}to be perplexed\textquotedblright{}; the native verb \emph{paḍu}, \textquotedblleft{}to feel, to suffer\textquotedblright{} (DEDR 3853) completes the verb.
\end{cousinweb}

\begin{grammarlens}[title={What you've built}]
One more word for describing what you lived and felt.
\end{grammarlens}

\subsection*{Your turn}
\begin{itemize}
\raggedright
  \item \emph{Say it:} \emph{tikamakapaḍu}
  \item \emph{Say it:} \emph{tikamakapaḍu}, once more
  \item \emph{Say it:} \emph{kotta nagaraṁlō dāri teliyaka nēnu tikamakapaḍḍānu.} — Not knowing the way in the new city, I got muddled.
  \item \emph{Recall:} say the Telugu for curiosity, then the Telugu for wits, then say \emph{tikamakapaḍu} again
\end{itemize}

\begin{tcolorbox}[breakable,colback=teal!4,colframe=teal!35!black,title={Before you move on}]
What does \te{తికమకపడు} mean? (\textquotedblleft{}To be confused; to be muddled\textquotedblright{}.) Say it once more.
\end{tcolorbox}
\section[\texorpdfstring{\te{తలచుకో} (talacukō)}{తలచుకో (talacukō)}]{\te{తలచుకో} (talacukō) — to think back on; to recall fondly}

\label{lesson:TE-C511-talacuko}



\begin{quote}
Before the new one: say the Telugu for wits, then the Telugu for to be confused.
\end{quote}

\subsection*{You'll want to know: \te{తలచుకో}}
\textbf{\te{తలచుకో}} — \emph{talacukō} — \textquotedblleft{}to think back on; to recall fondly\textquotedblright{}.

\begin{quote}
\te{చిన్నప్పటి} \te{రోజులను} \te{తలచుకుని} \te{అమ్మమ్మ} \te{నవ్వింది}.
\end{quote}

\begin{quote}
\emph{cinnappaṭi rōjulanu talacukuni ammamma navvindi.}
\end{quote}

\begin{quote}
Thinking back on her childhood, grandma smiled.
\end{quote}

Past (I …): \te{తలచుకున్నాను}, \emph{talacukunnānu}.

Often -\te{ని} \te{తలచుకుని}, thinking of: fondly, or with regret.

\begin{cousinweb}
\textbf{\te{తలచుకో}} (\emph{talacukō}) is built on \te{తలచు}, \textquotedblleft{}to think of\textquotedblright{}, with the reflexive -\te{కో}; the origin of \te{తలచు} is not settled.
\end{cousinweb}

\begin{grammarlens}[title={What you've built}]
One more word for describing what you lived and felt.
\end{grammarlens}

\subsection*{Your turn}
\begin{itemize}
\raggedright
  \item \emph{Say it:} \emph{talacukō}
  \item \emph{Say it:} \emph{talacukō}, once more
  \item \emph{Say it:} \emph{cinnappaṭi rōjulanu talacukuni ammamma navvindi.} — Thinking back on her childhood, grandma smiled.
  \item \emph{Recall:} say the Telugu for wits, then the Telugu for to be confused, then say \emph{talacukō} again
\end{itemize}

\begin{tcolorbox}[breakable,colback=teal!4,colframe=teal!35!black,title={Before you move on}]
What does \te{తలచుకో} mean? (\textquotedblleft{}To think back on; to recall fondly\textquotedblright{}.) Say it once more.
\end{tcolorbox}
\section[\texorpdfstring{\te{నెమరువేయు} (nemaruvēyu)}{నెమరువేయు (nemaruvēyu)}]{\te{నెమరువేయు} (nemaruvēyu) — to go over again; to chew the cud}

\label{lesson:TE-C511-nemaruveyu}



\begin{quote}
Before the new one: say the Telugu for to be confused, then the Telugu for to think back on.
\end{quote}

\subsection*{You'll want to know: \te{నెమరువేయు}}
\textbf{\te{నెమరువేయు}} — \emph{nemaruvēyu} — \textquotedblleft{}to go over again; to chew the cud\textquotedblright{}.

\begin{quote}
\te{ఆ} \te{సంతోషకరమైన} \te{రోజును} \te{నేను} \te{మళ్ళీ} \te{మళ్ళీ} \te{నెమరువేశాను}.
\end{quote}

\begin{quote}
\emph{ā santōṣakaramaina rōjunu nēnu maḷḷī maḷḷī nemaruvēśānu.}
\end{quote}

\begin{quote}
I went over that happy day again and again.
\end{quote}

Past (I …): \te{నెమరువేశాను}, \emph{nemaruvēśānu}.

First said of a cow chewing the cud, then of turning a memory over and over.

\begin{cousinweb}
\textbf{\te{నెమరువేయు}} (\emph{nemaruvēyu}) is native Dravidian: DEDR 3595 lists Telugu \emph{nemaru}, \textquotedblleft{}rumination, chewing the cud\textquotedblright{}, beside \emph{namalu}, \textquotedblleft{}to chew\textquotedblright{}; the verb \emph{vēyu}, \textquotedblleft{}to put\textquotedblright{}, makes it an action.
\end{cousinweb}

\begin{grammarlens}[title={What you've built}]
One more word for describing what you lived and felt.
\end{grammarlens}

\subsection*{Your turn}
\begin{itemize}
\raggedright
  \item \emph{Say it:} \emph{nemaruvēyu}
  \item \emph{Say it:} \emph{nemaruvēyu}, once more
  \item \emph{Say it:} \emph{ā santōṣakaramaina rōjunu nēnu maḷḷī maḷḷī nemaruvēśānu.} — I went over that happy day again and again.
  \item \emph{Recall:} say the Telugu for to be confused, then the Telugu for to think back on, then say \emph{nemaruvēyu} again
\end{itemize}

\begin{tcolorbox}[breakable,colback=teal!4,colframe=teal!35!black,title={Before you move on}]
What does \te{నెమరువేయు} mean? (\textquotedblleft{}To go over again; to chew the cud\textquotedblright{}.) Say it once more.
\end{tcolorbox}
\section[\texorpdfstring{\te{గుర్తుచేయు} (gurtucēyu)}{గుర్తుచేయు (gurtucēyu)}]{\te{గుర్తుచేయు} (gurtucēyu) — to remind}

\label{lesson:TE-C511-gurtuceyu}



\begin{quote}
Before the new one: say the Telugu for to think back on, then the Telugu for to go over again.
\end{quote}

\subsection*{You'll want to know: \te{గుర్తుచేయు}}
\textbf{\te{గుర్తుచేయు}} — \emph{gurtucēyu} — \textquotedblleft{}to remind\textquotedblright{}.

\begin{quote}
\te{పుట్టినరోజు} \te{సంగతి} \te{చెల్లెలు} \te{నాకు} \te{గుర్తుచేసింది}.
\end{quote}

\begin{quote}
\emph{puṭṭinarōju saṅgati cellelu nāku gurtucēsindi.}
\end{quote}

\begin{quote}
My younger sister reminded me about the birthday.
\end{quote}

Past (I …): \te{గుర్తుచేశాను}, \emph{gurtucēśānu}.

Who is reminded takes -\te{కి}: \te{నాకు} \te{గుర్తుచేసింది}. \te{గుర్తు} is a mark or a memory.

\begin{cousinweb}
\textbf{\te{గుర్తుచేయు}} (\emph{gurtucēyu}) is native Dravidian: DEDR 1847 lists Telugu \emph{guṟutu}, \textquotedblleft{}a mark, a sign\textquotedblright{}, and \emph{gurt(u) uṇḍu}, \textquotedblleft{}to be remembered\textquotedblright{}, with Tamil \emph{kuṟi}; \emph{cēyu}, \textquotedblleft{}to do\textquotedblright{}, makes the verb.
\end{cousinweb}

\begin{grammarlens}[title={What you've built}]
One more word for describing what you lived and felt.
\end{grammarlens}

\subsection*{Your turn}
\begin{itemize}
\raggedright
  \item \emph{Say it:} \emph{gurtucēyu}
  \item \emph{Say it:} \emph{gurtucēyu}, once more
  \item \emph{Say it:} \emph{puṭṭinarōju saṅgati cellelu nāku gurtucēsindi.} — My younger sister reminded me about the birthday.
  \item \emph{Recall:} say the Telugu for to think back on, then the Telugu for to go over again, then say \emph{gurtucēyu} again
\end{itemize}

\begin{tcolorbox}[breakable,colback=teal!4,colframe=teal!35!black,title={Before you move on}]
What does \te{గుర్తుచేయు} mean? (\textquotedblleft{}To remind\textquotedblright{}.) Say it once more.
\end{tcolorbox}
